\chapter{INTRODUCTION}

\section{Background}

The El Ni\~{n}o-Southern Oscillation (ENSO) is the dominant mode of interannual climate variability, influencing weather patterns worldwide through coupled ocean-atmosphere interactions in the tropical Pacific \cite{ham2019deep}. ENSO cycles between warm (El Ni\~{n}o) and cool (La Ni\~{n}a) phases, driven primarily by anomalies in Sea Surface Temperature (SST) across the equatorial Pacific Ocean. These events have far-reaching teleconnections that affect precipitation, temperature, and extreme weather events across multiple continents.

Furthermore, anthropogenic climate change is increasingly modulating ENSO dynamics. Recent climatological studies indicate a trend toward more frequent and intense extreme ENSO events, including a higher occurrence of Central Pacific (Modoki) El Ni\~{n}os and ``Super El Ni\~{n}os'' \cite{cai2015continued}. This non-stationarity in the background state of the tropical Pacific complicates the physical mechanisms driving ENSO, thereby exacerbating the ``spring predictability barrier'' and challenging the skill of traditional dynamical models.

South Asia, where the Indian Summer Monsoon provides over 70\% of annual rainfall, is highly vulnerable to ENSO-induced climate anomalies \cite{ye2024resonet}. El Ni\~{n}o events are associated with weakened monsoon circulation, leading to droughts in the Indian subcontinent, while La Ni\~{n}a events often intensify monsoon rainfall \cite{mahesh2024physics}. Climate outlooks indicate the possibility of ENSO-related anomalies during 2026, highlighting the importance of improved forecasting systems.

Nepal, situated in the central Himalayas, is particularly susceptible to these teleconnections due to its fragile topography and heavy reliance on monsoon-dependent agriculture and hydropower. ENSO phases significantly dictate the spatial and temporal distribution of precipitation across Nepal's diverse elevational gradients. During El Ni\~{n}o years, Nepal frequently experiences below-normal monsoon rainfall and delayed monsoon onsets, which critically deplete soil moisture, reduce hydropower generation capacity, and trigger winter and spring droughts. Conversely, La Ni\~{n}a phases are often associated with above-normal rainfall, increasing the risk of catastrophic flooding and landslides in the steep river basins \cite{shrestha2021enso}.

Traditional dynamical climate models, such as General Circulation Models (GCMs), provide seasonal ENSO forecasts but face limitations in prediction skill beyond 6--9 months due to the ``spring predictability barrier'' \cite{dijkstra2019deep}. Recent advances in deep learning have demonstrated superior forecasting skill at longer lead times (12--22 months) by extracting complex spatiotemporal patterns from vast reanalysis archives \cite{ham2019deep, mahesh2024physics}. These data-driven approaches leverage Convolutional Neural Networks (CNNs) and Long Short-Term Memory (LSTM) networks to capture nonlinear dynamics that traditional statistical methods and dynamical models struggle to represent.

\section{Motivation}
The motivation for this project stems from three critical factors:

\begin{itemize}
    \item \textbf{Imminent Climate Risk:} The 2026 El Ni\~{n}o poses significant risks to South Asia's agriculture, water resources, and disaster management infrastructure. Early and accurate prediction enables proactive mitigation strategies.
    
    \item \textbf{Limitations of Existing Models:} While dynamical models provide reasonable short-term ENSO forecasts, their skill degrades significantly beyond one season. Deep learning models have shown promise in extending this predictability horizon but are often not tailored to regional impact assessment.
    
    \item \textbf{Need for Actionable Regional Forecasts:} Most existing ENSO prediction systems focus on the Ni\~{n}o 3.4 index alone without providing spatially resolved impact assessments for vulnerable regions. There is a critical gap between global ENSO state prediction and localized climate impact forecasting.
\end{itemize}

\section{Problem Statement}
Despite significant progress in ENSO forecasting using deep 
learning, existing approaches face key limitations:

\begin{itemize}
    \item Most models predict only the Ni\~{n}o 3.4 index 
    without translating predictions into regional climate 
    impacts, limiting their practical utility for 
    disaster preparedness and agricultural planning.
    
    \item No unified deep learning framework currently exists 
    that couples ENSO state forecasting with deterministic 
    regional impact assessment specifically for South Asian 
    precipitation and temperature anomalies.
\end{itemize}

This project addresses these gaps by developing a multi-task 
CNN-TCN framework that forecasts ENSO state at 6 month lead times 
and simultaneously produces deterministic precipitation and 
temperature anomaly maps over South Asia.

\section{Objectives}
The primary objectives of this project are as follows:
\begin{itemize}
    \item To develop a CNN-TCN model to forecast the 
    Ni\~{n}o 3.4 index at a 6-month lead time using 
    ERA5 and ORAS5 reanalysis data spanning 1980--2025.
    
    \item To develop a coupled impact assessment module 
    that produces deterministic forecasts of 
    June--July--August--September precipitation 
    and temperature anomalies over South Asia based on 
    the predicted Ni\~{n}o 3.4 index.
\end{itemize}

\section{Scope and Limitations}

\subsection{Scope}
\begin{itemize}
    \item The system will forecast ENSO states using the Ni\~{n}o 3.4 index at lead times of 6 months.
    \item Regional impact assessment will cover the South Asian domain (5$^{\circ}$N--40$^{\circ}$N, 60$^{\circ}$E--100$^{\circ}$E).
    \item The model will be trained on publicly available datasets including ERA5 and ORAS5.
    \item The final deliverable will include a functional forecasting pipeline implemented in Python.
\end{itemize}

\subsection{Limitations}
\begin{itemize}
    \item The model is trained on historical data and may not capture unprecedented climate dynamics.
    \item The system does not account for other climate drivers (e.g., Indian Ocean Dipole, Arctic Oscillation) that may modulate ENSO impacts.
    \item Real-time operational deployment is beyond the scope of this project.
\end{itemize}

\section{Project Applications}
\begin{itemize}
    \item \textbf{Agricultural Planning:} The system provides early warning for monsoon deficiency or excess rainfall, enabling farmers and agricultural agencies to guide crop planning and irrigation management.

    \item \textbf{Disaster Preparedness:} The framework identifies regions at elevated flood or drought risk, allowing disaster management authorities to undertake proactive resource allocation.

    \item \textbf{Water Resource Management:} Seasonal reservoir operation planning is supported through deterministic precipitation forecasts generated by this system.

    \item \textbf{Climate Research:} The project contributes to the broader scientific understanding of ENSO teleconnections and advances AI-based climate prediction methodologies.
\end{itemize}
